Собственно sainte\_lague.cpp и является причиной, почему я выбрал именно эту тему для задания. Я положил его в репозиторий по приколу, чтобы был, и так он и лежал в корне один, никуда не принадлежащий, ни рыба ни мясо. Вот я и решил найти ему применение и перенести в какую-нибудь директорию.

Что касается самого кода, выглядит он довольно ужасно. Как понятно из названия, изначально программа предполагалась для расчета только метода Сент-Лагю, но потом я решил, что для сравнения лучше сразу выводить два метода рядом. Отсюда и дублированние переменные по типу \texttt{seat\_num1},\texttt{extra\_votes1} и т.д.

\inputminted[fontsize=\small, breaklines=true, style=borland, linenos]{cpp}{sainte_lague.cpp}